\chapter{Uso responsable de herramientas de inteligencia artificial}
\label{anexo:uso_ia}

Durante el desarrollo de este trabajo, se han utilizado herramientas de inteligencia artificial generativa como apoyo en distintas tareas. Una de ellas ha sido la búsqueda de referencias útiles, ya que ChatGPT no solo permite buscar información, sino también localizar fuentes a las que acudir para consultarla.

Además, estas herramientas han resultado realmente útiles durante el desarrollo en Unity, ya que permiten obtener de forma rápida explicaciones sobre el funcionamiento del motor y resolver dudas que, de otro modo, habrían requerido depender de foros, vídeos o tutoriales. También se han utilizado para obtener ideas sobre el diseño y la distribución de los elementos de la interfaz. Este apoyo ha sido fundamental en la elección de colores debido al daltonismo del autor.

\par\medskip
\textbf{Nota: Queda un poco raro decir daltonismo del autor, pero no puedo hablar en primera persona de mí; no sé como decirlo de otra forma.}
\par\medskip

La memoria se ha redactado en Overleaf, que incluye una herramienta de corrección automática mediante inteligencia artificial. Esta función ayuda a detectar problemas de puntuación, especialmente relacionados con el uso del punto y coma, y propone cambios para mejorar la formalidad y la claridad de la redacción, de una forma similar a los correctores ortográficos habituales de Word.

También se han utilizado estas herramientas para comprender conceptos nuevos relacionados con las distintas áreas del proyecto. Algunos artículos presentaban una dificultad elevada, y mantener una conversación con una inteligencia artificial generativa ha permitido comprender mejor determinados conceptos mediante explicaciones más sencillas y adaptadas a cada duda.

En ningún caso la inteligencia artificial sustituye el trabajo del ingeniero, ya que las decisiones, la implementación y la comprobación de los resultados siguen siendo responsabilidad del autor. Sin embargo, utilizada de forma crítica y responsable, puede ser una herramienta de gran utilidad para adquirir conocimientos.


\par\medskip
\textbf{Nota: Fui a ver las presentaciones de los TFGs de informática de julio, y el tribunal les preguntó directamente cómo habían usado la IA en el trabajo, ya que en la memoria no lo aclaraban. No sé si esto queda demasiado forzado, ya que algunos profesores son anti IA, pero como lo hablamos en una reunión, lo dejo aquí.}
\par\medskip